% General definitions and counters. Concrete graph and protocol: experimental-appendix.tex.
\section{Intervention definitions}
\label{app:interventions}

\subsection{Thresholding nonlinearities}

For finite $\kappa\geq0$, both thresholding nonlinearities retain equality and leave retained values unchanged.
Thus $G_{+,\kappa}$ is not $\operatorname{ReLU}(x-\kappa)$; at positive
cutoff both maps are discontinuous at their boundaries. The reported A4/A7
recipes use the one-sided map at $a,m,h,z$, including the attention
projection sites $a,z$, and the symmetric map only at $q,k,v$ in A7.
The comparison mask is detached: the derivative with respect to the input
is one for retained entries and zero for rejected entries. Consequently,
$G_{+,0}$ and ReLU have identical forward values but differ at zero, where
their derivatives are one and zero, respectively. $G_{\pm,0}$ is exactly the
identity in both values and gradients. These nonlinearities use neither a straight-through
estimator nor a top-$k$ selection rule. Unselected sites retain their original transformations.

\subsection{Activation pressure and OL1}
\label{app:pressure}

For the $J$ targeted site--layer tensors $A_j$, containing $n_j$ entries each,
the pressure objective is
$\mathcal{L}_1=J^{-1}\sum_{j=1}^{J}n_j^{-1}\sum_{i=1}^{n_j}|A_{j,i}|$.
The tensors contain the executed site activations, after the nonlinearity where one is applied.
Absolute values and means are evaluated in float32.
Each tensor receives equal weight regardless of its size. Adding pressure
sites changes the aggregate direction and, when the cap is inactive, its scale.
Task and pressure losses are averaged over the same accumulation microbatches.
For naive L1, the accumulated gradient of
$\mathcal{L}_{\mathrm{task}}+\lambda\mathcal{L}_1$ is globally norm-clipped
before AdamW~\citep{loshchilov2019adamw}.

For OL1, let $g_t$ and $g_1$ be the accumulated task and unweighted pressure
gradients. We first clip $g_t$ by its global norm and perform AdamW using only
that clipped task gradient. Let $\widehat m_t$ and $\widehat v_t$ be its newly
updated, bias-corrected moments. On parameters having both gradients and
initialized optimizer state, define the adaptive task and pressure vectors
$u=\widehat m_t/(\sqrt{\widehat v_t}+\epsilon_A)$ and
$w=g_1/(\sqrt{\widehat v_t}+\epsilon_A)$. The pressure gradient is unclipped
and has no first-moment buffer. Inner products and norms below pool all these
eligible parameters. With $c=\langle u,w\rangle$, $q=\|u\|_2^2$, and numerical
stabilizer $\epsilon>0$, the conflict-conditioned projection is
\begin{equation}
\widetilde w=
\begin{cases}
w-\dfrac{c}{q+\epsilon}u, & c<0\ \text{and}\ q>\epsilon,\\[3pt]
w, & \text{otherwise}.
\end{cases}
\label{eq:ol1-projection}
\end{equation}
For pressure weight $\lambda\geq0$ and required budget $b>0$, define the norm ratio
$r=\lambda\|\widetilde w\|_2/(\|u\|_2+\epsilon)$. Set
$s=\min\{1,b/(r+\epsilon)\}$ when $r>0$, or $s=1$ otherwise. After AdamW,
each eligible parameter group $p$, with learning rate $\alpha_p$, receives
$\theta_p\leftarrow\theta_p-\alpha_p\lambda s\widetilde w_p$.
The experiments' clipping thresholds, $b$, and $\lambda$ are condition
parameters.

The executed FP16 boundaries accumulate and unscale the two gradients
separately. A nonfinite gradient skips the entire boundary, including the
correction. Missing gradients or uninitialized moment state exclude a parameter
from its pooled geometry. All reported OL1 runs use $b=1$ and
$\epsilon=10^{-12}$; the local study sweeps $\lambda\in\{0.05,0.1,0.5,1\}$,
whereas multisite runs fix $\lambda=1$. The norms precede group learning rates
and exclude AdamW's decoupled weight decay.

\paragraph{Protected direction and saturation.}
For the ideal rule, omit stabilizers and take $u\ne0$. Project only under
conflict, using $\widetilde w=w-(u^\top w/\|u\|^2)u$; otherwise retain $w$.
Then the pressure descent vector is
\begin{equation}
d_{\mathrm p}=\min\!\left(\lambda,\frac{b\|u\|}{\|\widetilde w\|}\right)
\widetilde w,
\qquad d_{\mathrm p}=0\ \text{if }\widetilde w=0.
\label{eq:ol1-ideal}
\end{equation}
It satisfies $u^\top d_{\mathrm p}\geq0$, with equality after a conflicting
projection, and hence $u^\top(u+d_{\mathrm p})\geq\|u\|^2$.
Also $\|d_{\mathrm p}\|\leq b\|u\|$: the auxiliary vector is capped,
while the combined vector can be longer than $u$. At a fixed gradient pair
and optimizer state, $\lambda\|\widetilde w\|\geq b\|u\|$ makes the vector
independent of further increases in $\lambda$. A common learning rate
$\alpha$ turns this into the added displacement $-\alpha d_{\mathrm p}$.

With $\rho=b\|u\|$ and nonzero $\widetilde w$, the saturated vector maximizes
$w^\top d$ subject to $u^\top d\geq0$ and $\|d\|\leq\rho$.
Without conflict this follows from Cauchy--Schwarz. Under conflict,
$w=\widetilde w+cu/\|u\|^2$ with $c<0$, so every feasible $d$ satisfies
$w^\top d\leq\widetilde w^\top d\leq\rho\|\widetilde w\|$;
the stated vector attains both bounds. If $w$ exactly opposes $u$, zero
is an optimum, possibly nonunique. This is alignment with preconditioned
pressure, not necessarily the first-order decrease $g_1^\top d$.

The protected adaptive direction generally differs from the current task
gradient: first-order task loss depends on $g_t^\top d_{\mathrm p}$, and
finite steps also depend on curvature. Thus the property does not guarantee
task-loss preservation. In the implementation, stabilizers make the projection
approximately orthogonal, and $q\leq\epsilon$ bypasses it altogether.

\paragraph{Observed cap activity and objective scaling.}
% Analysis 020: 01_audit_logs.py; observations/O001-optimizer-logs.md.
All 34 OL1 conditions retain all 712 boundary records, with no skipped updates.
The cap is active exactly when $r>0$ and $b/(r+\epsilon)<1$, using the logged
pre-learning-rate ratio on eligible parameters. It binds on
7,283 of 24,208 boundaries (30.09\%), but this pooled fraction hides large
recipe differences. Across thresholds, A7-OL1 binds on 84.97--100\% of steps
at 14M and 100\% at 70M, versus 0.42--0.84\% at 410M. A4-OL1 binds on
0.14--2.67\% across all sizes; local h-only OL1 ranges from 0 to 10.67\%
across its four weights. The first boundary has zero learning rate: its
cap diagnostic describes a vector, not parameter movement. Supplementary
records retain counts, correction-ratio distributions and separate early
(steps 1--142) and late (570--712) summaries for every condition.

Multiplying a fixed objective by a positive scalar preserves the saturated
ideal vector's direction and norm when both versions remain saturated.
For a fixed seven-site target set, replacing $\sum_7/(7L)$ by
$\sum_7/(4L)$ only rescales that vector before the cap. Expanding four
targets to seven instead changes their gradient mixture, even under
saturation. Smaller coefficients in the written mean therefore do not
establish proportional dilution of the effective update. Relative weights
between old and new targets would test a different direction; the current
experiments do not identify that contribution separately.

\section{Sparsity accounting}
\label{app:accounting}

\subsection{Exact counters}

Activation sparsity at site $s$ is
$\mathcal{Z}_s=\sum_{\ell,e,i}\mathbf{1}\{A^{(\ell,e)}_{s,i}=0\}/
\sum_{\ell,e}|A_s^{(\ell,e)}|$, where $\ell$ indexes layers, $e$ evaluation
batches, $i$ tensor entries, and $|A|$ the number of elements in tensor $A$. Near-zero mass replaces equality by
$|A_{s,i}|\leq\varepsilon$ for a stated tolerance; it does not contribute to
exact-zero counts unless the value is zero. All aggregate fractions pool
integer numerators and denominators across batches and layers before division.

The atomic unit of model-wide accounting is one scalar multiplication in a
matrix product. If $X\in\mathbb{R}^{n\times p}$ feeds a linear projection with
$q$ outputs, the count is $N=npq$, with $N_0=q\sum_{i,j}\mathbf{1}\{X_{ij}=0\}$.
Zero weights receive no credit. For the multihead attention graph specified in Appendix~\ref{app:sites},
both operands are activations,
so a product is credited once if either is zero. Within one block and
evaluation batch, for query position $t$,
key position $u\leq t$, and within-head coordinate $j$,
\begin{align}
N_{0,QK}&=\sum_{b,h,t}\sum_{u\leq t}\sum_j
\mathbf{1}\{Q_{bhtj}=0\ \lor\ K_{bhuj}=0\},\label{eq:qk-count}\\
N_{0,PV}&=\sum_{b,h,t}\sum_{u\leq t}\sum_j
\mathbf{1}\{P_{bhtu}=0\ \lor\ V_{bhuj}=0\}.\label{eq:pv-count}
\end{align}
Here $P$ is the causal softmax probability tensor; the model numerator sums
these counts over blocks and evaluation batches. Future masked positions
are excluded from both counts, while zeros on valid positions, including
floating-point underflow, are included. Counts use the operands actually consumed by each matrix product; upstream
masks do not determine those operands' zeros after intervening transformations.
Earlier value positions participate
in more causal pairs, so local value sparsity alone does not determine $PV$
sparsity.

\subsection{Architectural reach and normalization}
\label{app:ceiling}

For a specified graph and workload, let $N_o$ count the valid scalar
multiplications in block operation family $o$, pooled over layers and
examples. Then $N_{\mathrm{model}}=\sum_o N_o+N_{\mathrm{head}}$, while $N_0$
sums the corresponding zero-operand counts only in block operations.
The block-only view is $\Sblock=N_0/\sum_o N_o$. Both sparsities are fractions;
percentages multiply them by 100. Embedding lookup, normalization,
nonlinearities, softmax, RoPE, bias and residual additions, and data movement
are outside this accounting.

To construct the reach numerator, count products guaranteed to have a zero
activation operand as a structural consequence of 100\% activation sparsity at the selected sites, independently of checkpoint values or natural zeros elsewhere.
Include exact zero-preserving propagation under the declared graph rules. Products reached through more
than one site receive one credit. The resulting count is
$N_{\mathrm{reach}}(\mathcal{N})$. Division by the same $N_{\mathrm{model}}$
gives $\Rarch$. Natural zeros outside this selected-site reach can contribute
to observed $\Smodel$, so architectural reach is not an upper bound on every
observation. Its gap to an observation does not measure sparsity available
without quality loss. Appendix~\ref{app:architecture-counts} specializes
the denominator and reach rules to the Pythia graph used in the experiments.
